%% ----------------------------------------------------------------------------
%%  §9.11  Punti fondamentali del capitolo
%% ----------------------------------------------------------------------------
\section{Punti fondamentali del capitolo}
\label{sec:sintesi-cap09}

Integrazione e derivazione sono operazioni inverse: la funzione integrale di una
funzione continua ne è una primitiva, e un integrale definito si calcola come differenza
dei valori di una primitiva negli estremi. Il calcolo degli integrali diventa così la
ricerca delle primitive, per la quale bastano poche tecniche generali, ciascuna figlia di
una regola di derivazione.

\begin{puntichiave}
\item \textbf{Funzione integrale.} Se $f$ è integrabile in $[a,b]$ e $\abs f\le K$, la
funzione $F(x)=\int_{a}^{x}f(t)\,dt$ è lipschitziana con costante $K$, quindi continua
(\cref{prop:integrale-lipschitz}); può però non essere derivabile: un salto di $f$
diventa un punto angoloso di $F$.
\item \textbf{Teorema fondamentale.} Se $f$ è continua in $[a,b]$, allora $F$ è
derivabile e $F'(x)=f(x)$ per ogni $x\in[a,b]$ (\cref{teo:fondamentale}); se $f$ è soltanto
integrabile, basta la continuità in un punto $x_{0}$ per avere $F'(x_{0})=f(x_{0})$
(\cref{oss:fondamentale-puntuale}). Con la regola della
catena, $\frac{d}{dx}\int_{c}^{\varphi(x)}f(t)\,dt=f(\varphi(x))\,\varphi'(x)$
(\cref{cor:estremi-variabili}).
\item \textbf{Primitive.} Ogni funzione continua in un intervallo ha primitive; in un
intervallo due primitive della stessa funzione differiscono per una costante
(\cref{teo:caratterizzazione-primitive}). Fuori dagli intervalli la tesi è falsa: su
$\R\setminus\set{0}$ la costante può cambiare tra $x<0$ e $x>0$.
\item \textbf{Formula fondamentale.} Se $f$ è continua in $[a,b]$ e $G'=f$, allora
$\int_{a}^{b}f(x)\,dx=G(b)-G(a)$ (\cref{teo:formula-fondamentale}); la formula vale anche
per $f$ soltanto integrabile, purché dotata di primitiva
(\cref{oss:formula-fondamentale-integrabili}).
\item \textbf{Integrale indefinito.} $\int f(x)\,dx=F(x)+c$ indica l'insieme delle
primitive in un intervallo (\cref{def:integrale-indefinito}): è un insieme di funzioni,
non un numero. Si calcola con la linearità, gli integrali immediati
(\cref{tab:integrali-immediati}) e la regola $\int g(\varphi(x))\,\varphi'(x)\,dx
=G(\varphi(x))+c$.
\item \textbf{Funzioni razionali.} Si divide il numeratore per il denominatore e si
scompone il resto in fratti semplici; con denominatore di secondo grado si ottengono
logaritmi se $\Delta>0$, un termine $-\frac{mx_{1}+n}{x-x_{1}}$ se $\Delta=0$, un logaritmo e
un'arcotangente se $\Delta<0$ (\cref{prop:razionali-secondo-grado}).
\item \textbf{Integrazione per parti.} Se $f$ e $g$ hanno derivata continua,
$\int fg'\,dx=fg-\int f'g\,dx$ (\cref{teo:per-parti}), e negli integrali definiti compare
il termine $\bigl[fg\bigr]_{a}^{b}$ (\cref{teo:per-parti-definito}). Come fattore finito
si sceglie una funzione che si semplifica derivando; talvolta ricompare l'integrale di
partenza, o uno dello stesso tipo (formule di riduzione).
\item \textbf{Integrazione per sostituzione.} Se $f$ è continua e $g$ ha derivata
continua, $F\circ g$ è una primitiva di $(f\circ g)\,g'$ (\cref{teo:sostituzione}); negli
integrali definiti $\int_{a}^{b}f(x)\,dx=\int_{c}^{d}f(g(t))\,g'(t)\,dt$ con $g(c)=a$,
$g(d)=b$, senza che $g$ debba essere invertibile (\cref{teo:sostituzione-definita}).
\item \textbf{Errori da evitare.} Una primitiva vale in un intervallo: $\int\frac{dx}x
=\log\abs x+c$, non $\log x+c$, e i programmi di calcolo simbolico omettono spesso il valore
assoluto. Non tutte le funzioni integrabili hanno primitive (la funzione che vale $-1$ per
$x<0$ e $1$ per $x\ge0$ non ne ha, per il lemma di Darboux), e le primitive di $e^{-x^{2}}$ o
di $\frac{\sen x}x$ esistono ma non sono elementari.
\item \textbf{Aree.} La regione compresa tra i grafici di $g\le f$ continue in $[a,b]$ ha
area $\int_{a}^{b}(f-g)\,dx$ \eqref{eq:area-tra-curve}; con le sostituzioni $x=\sen t$ e
$x=rt$ si ottengono l'area $\pi r^{2}$ del cerchio e l'area $\pi ab$ dell'ellisse
(\cref{es:cerchio-ellisse}).
\item \textbf{Taglio e momento flettente.} Per un carico $l$ continuo sulla semiala,
$T'=-l$, $M'=-T$, $T(s)=M(s)=0$ e $M(0)=\int_{0}^{s}y\,l(y)\,dy$
(\cref{prop:taglio-momento}); la distribuzione ellittica dà $M(0)=\frac{Wb}{3\pi}$, circa il
$15\%$ in meno del carico uniforme (\cref{es:taglio-ellittico}). Integrando si ricavano
anche il tempo di salita e la legge della pressione nella troposfera standard
(\cref{es:tempo-salita,es:pressione-isa}).
\item \textbf{Integrali impropri.} Per $f$ continua e non negativa in $[a,+\infty)$,
l'integrale improprio $\int_{a}^{+\infty}f$, limite di $\int_{a}^{b}f$ per $b\to+\infty$, esiste
sempre, finito o $+\infty$ (\cref{def:improprio-illimitato}); $\int_{1}^{+\infty}x^{-p}\,dx$ converge se e solo se
$p>1$, $\int_{0}^{1}x^{-p}\,dx$ se e solo se $p<1$
(\cref{es:impropri-p-infinito,es:impropri-p-zero}), e il criterio del confronto
(\cref{teo:confronto-impropri}) prova, ad esempio, la convergenza di
$\int_{0}^{+\infty}e^{-x^{2}}dx$.
\item \textbf{Il logaritmo come integrale.} Posto $L(x)=\int_{1}^{x}\frac{dt}t$ per $x>0$,
si ha $L'=\frac1x$, $L(x_{1}x_{2})=L(x_{1})+L(x_{2})$, $L(e)=1$, e $L$ è una biiezione di
$(0,+\infty)$ su $\R$ (\cref{prop:L-proprieta}); la sua inversa $E$ verifica $E'=E$
(\cref{prop:E-proprieta}), e $a^{x}=E(x\,L(a))$ ritrova la potenza a esponente reale
(\cref{def:potenza-integrale}).
\end{puntichiave}

\paragraph{Il filo del discorso.} Il calcolo integrale ci ha dato un modo nuovo di
rappresentare una funzione di classe $C^{1}$: come integrale della sua derivata,
$f(x)=f(x_{0})+\int_{x_{0}}^{x}f'(t)\,dt$. Nel \cref{cap:taylor} integreremo ripetutamente
per parti questa identità, ottenendo la formula di Taylor con il resto integrale e con il
resto di Lagrange, che permettono di \emph{stimare} l'errore delle approssimazioni
polinomiali; la stessa formula, applicata al momento flettente di questo capitolo, darà
la freccia di un'ala a sbalzo.
